\section{Problem and scope}
Let $D=\{x\in\mathbb R^2:\|x\|\le 1\}$. We define
\[
r_D(100)=\inf\left\{r:\ D\subseteq\bigcup_{i=1}^{100}B(c_i,r),\quad c_i\in\mathbb R^2\right\}.
\]
The problem is continuous and non-convex. The centres are unrestricted, and the combinatorial type of the induced Voronoi or covering decomposition changes during optimization. Consequently, a numerical configuration supplies an upper bound, whereas a lower bound must exclude every admissible configuration at a prescribed radius.

This work is a certificate-oriented computational study rather than a claim of a formalized theorem in a proof assistant. The final decision layer uses integer, rational, or outward interval arithmetic whenever possible. Floating-point optimization and high-precision numerics are discovery tools only.

\subsection{What is proved and what is not}
The public release certifies
\[
  0.113235<r_D(100)\le 0.1154469548855829.
\]
The lower inequality is a six-branch exclusion certificate at $R=22647/200000$. The upper inequality is a replayed explicit configuration. Neither statement identifies the exact optimum. In particular, the nearby target $0.11324$ remains open in the present certificate system.

The distinction matters: a successful local box exclusion is not a global proof; an unblocked numerical point is not a certificate; and a covering configuration does not prove optimality.
